\documentclass{article}
\title{Developers}
\author{Charles Lin}

\begin{document}

    \maketitle

    \section{Repository}
    \subsection{This repository contain three directory:}
    \begin{enumerate}
        \item Documentation This include all the documentation for the project.
        \\Include a README file for general question, and a Developer file introduce the repository and the purpose of each file.
        \item Grading This directory reserve for grading purpose.
        \item Source This directory is where the source code located, also include an Makefile for build and clean.
    \end{enumerate}

    \section{Phase 1}
    \subsection{Main.java}
    It contains the main, act as the entry for the program.\\
    simple read the argument and print properly output.
    \section{Phase 2}
    \subsection{utilties.Token.java}
    This Token class define an instance that contain four variable:
    \begin{enumerate}
        \item fileName which the original file that generate the token.
        \item lineNumber the original line that token located.
        \item tokenType the type of the token.
        \item lexeme the lexeme of the token.
    \end{enumerate}
    By override toString() function of this instance, allow me easy to format properly output.
    \\Also reserve for later stage of the project if need store the token as list.
    \subsection{Lexer.java}
    The Lexer class is the main feature implemented for phase 1.\\
    This lexer class will take input for a file name, scan and tokenize it.\\ \\
    \begin{enumerate}
        \item TOKEN\_MAP This hashMap contain a string as key which represent the
        two characters operator, keyword, or type; and integer type value that correspond of each key.
        Allow me quick look up the type of each key.
        \item static boolean encounterError and static string errorMessage are shared across the lexer class.
        when encounter an error, allow lexer throw out exception, pass the proper error message, and trace back if it was recursively called
        for include statement.
        \item The scanFile() will initialize the process of scan the file.
        \\It will catch if the encounterError is flagged. Once it detect the flag,
        it will throw the exception and handle it.\\
        This function control the flow of scan and tokenize the file.
        \item The processLine() will called by scanFile for each line.\\
        It will scan through the line, match each token by use regex expression.\\
        And check if there is any error, once detect the error it will flagged encounterError and return the control back to scanFile().\\
        If there is no error, it will generate output for each token either in output stream or output file.
    \end{enumerate}
    \subsection{Main.java}
    In the phase 2, it update the input that will take, when in the mode 1 and supply with a file, it will create a lexer and start scan.

    \section{Phase 3}
    \subsection{CFG}
    This is the context free program that I use to build my parser, it might not unambiguous, but the issue will handle by the code.\\
    $Program \rightarrow Declaration Program $\\
    $| FunctionPrototype Program$\\
    $| FunctionDefinition Program$\\
    $| StructDefinition Program$\\
    $| \varepsilon$\\
    $ListOfDeclaration \rightarrow Declaration | Declaration ListOfDeclaration$\\
    $Declaration \rightarrow Type ListOfIdentifierDeclaration ;$\\
    $FunctionPrototype \rightarrow Type Identifier ( ) ;$\\
    $| Type Identifier ( ListOfParameter ) ;$\\
    $FunctionDefinition \rightarrow Type Identifier ( ) \{ BlockStatement \}$ \\
    $| Type Identifier ( ListOfParameter ) \{ BlockStatement \}$\\
    $StructDefinition \rightarrow struct Identifier \{ \} ;$\\
    $| struct Identifier \{ ListOfDeclaration \} ;$\\
    $ListOfIdentifierDeclaration \rightarrow IdentifierDeclaration $\\
    $| IdentifierDeclaration , ListOfIdentifierDeclaration$\\
    $IdentifierDeclaration \rightarrow Identifier | Identifier [ Integer ]$\\
    $| Identifier = Expression | Identifier [ Integer ] = Expression$\\
    $Type \rightarrow void | char | int | float | struct$ \\//Note, the const can appear either before or after type keyword\\
    $ListOfParameter \rightarrow Parameter | Parameter , ListOfParameter$\\
    $Parameter \rightarrow Type Identifier | Type Identifier [ ]$\\
    $BlockStatement \rightarrow ListOfDeclaration BlockStatement| ListOfStatement BlockStatement | \varepsilon$\\
    $ListOfStatement \rightarrow Statement ListOfStatement | Statement$ \\
    $Statement \rightarrow ;$\\
    $| Expression ;$\\
    $| break ;$\\
    $| continue ;$\\
    $| return ; | return Expression ;$\\
    $| if ( Expression ) Statement | if ( Expression ) Statement else Statement$\\
    $| for ( Expression ; Expression ; Expression )$ //Note all Expression right here are optional.\\
    $| while ( Expression ) Statement$\\
    $| do Statement while ( Expression ) ;$\\
    $| \{ BlockStatement \}$\\
    $ArgumentList \rightarrow Expression , ArgumentList | Expression$\\
    $l-value \rightarrow Identifier . Identifier l-value | Identifier . Identifier [ Expression ]$\\
    $AssignmentOperator \rightarrow = | += | -= | *= | /=$\\
    $PostfixOperator \rightarrow ++ | -- | . Identifier | [ Expression ] | ( ArgumentList )$\\
    $RelationalOperator \rightarrow < | > | <= | >=$\\
    $AdditiveOperator \rightarrow + | -$\\
    $MultiplicativeOperator \rightarrow * | / | \%$\\
    $UnaryOperator \rightarrow ! | ~ | -$
    $Expression \rightarrow AssignmentExpression$\\
    $AssignmentExpression \rightarrow l-value AssignmentOperator Expression | conditionalExpression$\\
    $ConditionalExpression \rightarrow LogicalOrExpression | LogicalOrExpression ? Expression : Expression$\\
    $LogicalOrExpression \rightarrow LogicalAndExpression (|| LogicalAndExpression)*$ \\
    $LogicalAndExpression \rightarrow BitwiseOrExpression (\&\& BitwiseOrExpression)*$\\
    $BitwiseOrExpression \rightarrow BitwiseAndExpression (| BitwiseAndExpression)*$\\
    $BitwiseAndExpression \rightarrow EqualityExpression (\& EqualityExpression)*$\\
    $EqualityExpression \rightarrow RelationalExpression == RelationalExpression | RelationalExpression != RelationalExpression$\\
    $RelationalExpression \rightarrow AdditiveExpression (RelationalOperator AdditiveExpression)*$\\
    $AdditiveExpression \rightarrow MultiplicativeExpression (AdditiveOperator MultiplicativeExpression)*$\\
    $MultiplicativeExpression \rightarrow UnaryExpression (MultiplicativeOperator UnaryExpression)*$\\
    $UnaryExpression \rightarrow UnaryOperator UnaryExpression | ( Type ) UnaryExpression | PostfixExpression$\\
    $PostFixExpression \rightarrow PrimaryExpression (PostfixOperator)*$\\
    $PrimaryExpression \rightarrow Identifier | Literal | ( Expression )$\\
    $Literal \rightarrow Integer | RealNumber | Char | String$\\

    \subsection{Parser.java}
    \begin{enumerate}
        \item isInsideFunction, isInsideStruct, and isInsideParameterList are three global variable that use to track current scope.\\
        \item a new parse obj take input a list of tokens(generate by lexer), and a file name(string).\\
        \item The parse will scan through the list of token to check syntax and produce the final output.\\
        \item match(String expect) consume token optionally.\\
        \item expect(String expect) either match or panic.\\
        \item panic() print the error and exit the program.\\
        \item parseProgram() start parse.\\
        \item printOutput() to generate output.\\
    \end{enumerate}

    \section{Phase 4}
    The TypeChecker.java is based on the Parser.java but with modified implemented rules that check the type.\\
    It used four hash maps:\\
    globalVariable, functions, prototypes, and globalStruct.\\
    And two stacks that will push a hashmap when we enter braces and pop when we leave the braces:\\
    localVariable, localStruct.
    To maintain a symbol table and track scope. \\
    Also include new utilities class, Variable, Function, and Struct for meta info.

    \section{Phase 5}
    introduce many emit function to generate java assembly code. The emit function will be called and use during the parse process while the compiler consume the input file.\\

    \section{Phase 6}
    Continue on Phase 5 add branch control, spilt the old parseStatementBlock parse function into decoupling version.\\
    TypeCheck.java is for Phase 5 and 6, TypeChecker2.java is for phase 4 trace back.\\
\end{document}